\section{Harmonic restriction and cycle moments}
\label{rec:sec-harmonic}

The density estimates already proved describe the joint routes of tracked
cards.  We now retain more of the representation carried by those cards.
For a diagram whose first row has length $n-k$, its first occurrence on
ordered tuples has exactly $k$ slots.  Summing out any one slot therefore
annihilates that representation.  This harmonic cancellation controls how
much of its level can disappear when the network splits into two halves.

The first goal is a contraction on the scale $k\log(en/k)$.  A separate
cycle-moment estimate will then control the dimension of the diagram below
the first row.  We give two transfers from that moment, one by clipping
kernel entries and one by averaging matrix coefficients.  Both retain the
lower-diagram dimension, whereas the final tuple-trace argument can sum all
representations without separating that dimension.

Let $T_n$ be one sweep on $n=2^d$ positions, and put $K_n=T_n^*T_n$.
Reversing the coordinate order conjugates this law to the law of
$T_nT_n^*$ by the permutation that reverses the binary coordinates.  Thus
all statements below also hold for that orientation, with the same
constants and with the starting tuple relabeled by this permutation.
At the one-position base, the empty product gives $T_1=K_1=I$.
If $n=2m$, then
\begin{equation}\label{rec:network-recursion}
 K_n=P(K_m\otimes K_m)P,
\end{equation}
where $P$ independently averages the $m$ switches between corresponding
positions of the two halves.  Thus $K_n$ is positive, self-adjoint, and a
contraction in every unitary representation.  Throughout this section
traces count dimension, and logarithms are natural unless a base is shown.
The regular trace of an operator $A$ means
$\operatorname{tr}_{\rm reg}A=\sum_{\lambda\vdash n}D_\lambda
\operatorname{tr}_{V_\lambda}A$.  On a tuple space it means the ordinary
trace of its matrix in the tuple basis; rescaling all basis vectors from
counting to probability norm does not change it.
For $\lambda=(n-k,\beta)\vdash n$, write $D_\lambda=\dim V_\lambda$,
$f_\beta=\dim V_\beta$, and
\[
 \Lambda_n(k)=k\log(en/k)\quad(k>0),\qquad\Lambda_n(0)=0,\qquad
 W_n(\lambda)=\|K_n|_{V_\lambda}\|_{\op}
              =\|T_n|_{V_\lambda}\|_{\op}^{2}.
\]

\begin{theorem}\label{rec:main-harmonic-cycle}
There are absolute constants $c,C,\zeta>0$, a number $\delta\in(0,1)$,
and an integer $p_0$ with the following properties.  For every dyadic $n$
and $\lambda=(n-k,\beta)\vdash n$ with $k\ge1$,
\begin{align}
 W_n(\lambda)&\le e^{-c\Lambda_n(k)},\label{rec:level-contraction}\\
 W_n(\lambda)&\le e^{C\Lambda_n(k)}f_\beta^{-\zeta}.
 \label{rec:tail-contraction}
\end{align}
Independently, for every $1\le k\le n$, let $X_{n,k}$ be the ordered injective $k$-tuples of positions, with
uniform probability measure, and let
$\mathcal K_{n,k}(x,y)=(n)_k\Pr_{K_n}(x\mapsto y)$ be its kernel density.
Then
\begin{align}
 \sup_x\E_y\mathcal K_{n,k}(x,y)^{1+\delta}
   &\le e^{C\Lambda_n(k)},\label{rec:tuple-moment}\\
 \operatorname{tr}_{X_{n,k}}K_n^{p_0}
   &\le e^{C\Lambda_n(k)}.\label{rec:tuple-flat-trace}
\end{align}
Consequently, after changing absolute constants,
\begin{align}
 \|T_n|_{V_\lambda}\|_{\op}
   &\le \exp[-c\{\Lambda_n(k)+\log f_\beta\}],
       \label{rec:combined-tail-norm}\\
 W_n(\lambda)&\le e^{-c_1\Lambda_n(k)}D_\lambda^{-c_2}
       \label{rec:combined-dimension-norm}
\end{align}
for absolute $c_1,c_2>0$.
\end{theorem}

\subsection{Harmonic tuple restriction}
We first describe the probability law on the two child representations
that arises from one maximizing vector.  The law is a spectral weighting,
not an independent random choice of Young diagrams.  The removal estimates
below hold uniformly over that weighting.

The position action of $S_n$ on $X_{n,k}$ commutes with the action of
$S_k$ on the tuple labels.  Pieri's rule~\cite{Sagan2001,Stanley1999}
gives the multiplicity-free injection decomposition
\cite[Section 3.2, equation (3.1)]{DukesIhringerLindzey2020}: one copy of
$V_\mu\otimes V_\xi$ precisely when $\mu/\xi$ is a horizontal strip
of size $n-k$.  In particular the position type $(n-k,\beta)$ has
label type $\beta$ at this smallest tuple size.  It does not occur on
$X_{n,k-1}$.  Consequently its functions are harmonic: summing over any
one tuple entry, with all the others fixed, gives zero.

Work in the position-$\lambda$ isotypic space of $\ell^2(X_{n,k})$,
which is $V_\lambda\otimes V_\beta$ under the commuting position and
slot-permutation actions.  Use counting norm here and take a unit
eigenvector $v$ attaining a
positive value of $W_n(\lambda)$.
Equation~\eqref{rec:network-recursion} implies $Pv=v$.  Decompose $v$
by the set of labels in each half and then by the two position types.
The squared component norms define a probability law.  Write $a,b$ for
the numbers of labels in the halves, $a+b=k$, and $r,s$ for their first-row
deficits.  Branching and the tuple realization give
\begin{equation}\label{rec:restriction-law}
 0\le r\le a,\quad 0\le s\le b,\quad r,s\ge k-m,
 \qquad W_n(\lambda)\le\E[W_m(\mu)W_m(\nu)].
\end{equation}
The last lower bounds follow from $\mu,\nu\subseteq\lambda$.
The decompositions by label sets commute with the half-position groups,
so they do not create additional types.

\begin{lemma}[Removal and splitting]\label{rec:removal-splitting}
For the component law above, $a$ is a mixture of
$t+\operatorname{Bin}(k-2t,1/2)$.  In particular
\begin{equation}\label{rec:split-tail}
 \Pr(|a-k/2|\ge z)\le2e^{-2z^2/k},\qquad
 \E e^{u(a-b)^2/k}\le1+Cu\quad(0\le u\le1/4).
\end{equation}
For every fixed $\eta>0$ and integer $l\ge1$,
\begin{equation}\label{rec:loss-binomial-general}
 \E\!\left[\mathbf1_{\{m-a-r\ge\eta m\}}
                   {a-r\choose l}\right]
       \le(C_\eta\sqrt{k/n})^l,
\end{equation}
and the same statement holds in the other half.  Thus, with
$\alpha=1/8$ and
\[
 L_1=(a-r)\mathbf1_{\{a\le3m/4,\ r<\alpha m\}},\qquad
 L_2=(b-s)\mathbf1_{\{b\le3m/4,\ s<\alpha m\}},
\]
we have $\E{L_i\choose l}\le(C\sqrt{k/n})^l$ and
\begin{equation}\label{rec:loss-mgf}
 \E e^{u\log(en/k)L_i}\le1+Cu
 \quad(0\le u\le u_0)
\end{equation}
for an absolute $u_0>0$.  A second useful form, with $\alpha=1/16$, is
\begin{equation}\label{rec:loss-tail}
 \Pr\big((a-r)\mathbf1_{\{r<\alpha m\}}\ge x\big)
       \le C e^{-cx\log(en/k)}\qquad(x\in\mathbb Z_{\ge1}).
\end{equation}
\end{lemma}

\begin{proof}
Randomly orient the crossing pairs after sampling a configuration with
probability $|v|^2$.  A doubly occupied pair contributes one label to each
side, and each singly occupied pair contributes an independent fair
choice.  This proves the mixture and the concentration bounds.

For the removal bounds it suffices to take $1\le l\le k$; the
binomial coefficients vanish when $l>k$.  Fix the first half $S_0$.
For $J\subseteq[k]$, let $D_J^{S_0}$
sum the coordinates with labels in $J$ over $S_0$, leaving all the
other labeled coordinates in place.  Its target has counting measure
on the ordered injective $(k-|J|)$-tuples.  We use the zero extension
$\widetilde v$ of $v$ to write these sums without collision exclusions.
For a transversal $S$ of the matched pairs define $D_J^S$ in the same
way, and put
\[
 Q_l=\mathbb E_S\sum_{|J|=l}\|D_J^S v\|_2^2.
\]
Here $S=gS_0$ for a uniformly chosen element $g$ of the matching
switch group.  Since $Pv=v$, every such $g$ fixes $v$ and
$D_J^{gS_0}v=U_{k-l}(g)D_J^{S_0}v$, where $U_{k-l}(g)$ permutes
the retained positions.  Thus
\[
 Q_l=\sum_{|J|=l}\|D_J^{S_0}v\|_2^2.
\]
We will bound this one quantity below using the fixed-half
representation decomposition and above using the random transversal.
Write $D_i=D_{\{i\}}^{S_0}$.  In a sector
with $a$ labels and position type $\mu=(m-r,\theta)$,
\begin{equation}\label{rec:removal-quadratic}
 \sum_iD_i^*D_i\succeq(a-r)(m-a-r+1)I.
\end{equation}
Indeed the operator on the left is the arrangement adjacency operator
plus $aI$; the position--label content identity
\cite[Lemma 3.1 and Proposition 3.3]{AraujoBratten2017} gives
$aI+\Omega_{\rm pos}-\Omega_{\rm lab}-{m-a\choose2}I$, where $\Omega$
is the transposition class sum.  In the following calculation,
$c(\rho)=\sum_{(i,j)\in\rho}(j-i)$ denotes its content scalar,
rather than the complementary positive Casimir scalar of Section~\ref{sec:casimir}.
If the label type is $\xi$, Pieri gives
$\xi\supseteq\theta$.  Put $h=a-r$.  The content formula and successive
addition of its $h$ extra boxes give
\[
 c(\mu)={m-r\choose2}+c(\theta)-r,
 \qquad c(\xi)\le c(\theta)+hr+{h\choose2}.
\]
Substitution proves~\eqref{rec:removal-quadratic}.
Iterating over $l$ distinct labels and dividing by the $l!$ possible
orders gives the lower bound
\[
 Q_l
 \ge(\eta m)^l\E\!\left[
   \mathbf1_{\{m-a-r\ge\eta m\}}{a-r\choose l}\right].
\]
For clarity, in a sector whose left-label set is $A$, a deletion
contributes only when $J\subseteq A$.  Afterward the left-label set is
$A\setminus J$, so the target still distinguishes the original sector
when $J$ is fixed.  Deletion also intertwines the two half-position
groups, so different position types remain orthogonal.  After $j$
deletions the surviving type still has deficit $r$, while the occupancy
is $a-j$.  Iterating~\eqref{rec:removal-quadratic} therefore supplies
$(a-r)_{l}(\eta m)^l$ on the indicated event.  The $l!$ orders of
each deleted label set give the displayed binomial coefficient.

For the upper bound, give the two positions of each pair opposite fair
signs $\epsilon(x)$, so $S=\{x:\epsilon(x)=1\}$.  For fixed retained
tuple $y$ and deleted labels $J$, harmonicity gives the exact formula
\[
 D_J^S v(y)=2^{-l}\sum_z\widetilde v(y,z)
                         \prod_{i\in J}\epsilon(z_i).
\]
Indeed expanding the product of half-indicators leaves this term only:
every term missing a sign has a zero sum in that coordinate.
On squaring and averaging, two deleted assignments can correlate only
if their sets of pair indices with odd occupancy agree.  Suppose that
set has $l-2t$ members.  A partner assignment is specified by at most
$\binom mt$ choices of doubly occupied pair sites, two endpoint
choices at each of the $l-2t$ odd sites, and $l!$ assignments of the
chosen positions to the deleted labels.  Hence the class has size at
most $l!2^{l-2t}\binom mt$.  Ignoring conflicts with the retained
tuple only increases this bound.

Cauchy--Schwarz in each such class bounds its squared signed sum by
the class size times the sum of $|\widetilde v(y,z)|^2$.  Now sum
over $y$ and $J$.  In a fixed full $k$-tuple there are at most
$\lfloor k/2\rfloor$ doubly occupied pair sites.  Choose $t$ of
them for the deleted double pairs and then choose the remaining
$l-2t$ deleted labels.  Thus there are at most
$\binom{\lfloor k/2\rfloor}{t}\binom{k}{l-2t}$ choices, with
overcounting harmless.  Since $\|v\|_2=1$, this proves
\begin{equation}\label{rec:removal-upper-sum}
 Q_l
 \le4^{-l}\sum_{0\le t\le l/2}
 l!2^{l-2t}{m\choose t}{\lfloor k/2\rfloor\choose t}{k\choose l-2t}
 \le(C\sqrt{mk})^l.
\end{equation}
The last step uses $k\le2m$ and
$l!/(t!^2(l-2t)!)\le3^l$: each summand is at most
$C^l m^t k^{l-t}\le C'^l(mk)^{l/2}$, and the number of summands
is absorbed into the absolute constant.  Comparing the two bounds on
$Q_l$ proves
\eqref{rec:loss-binomial-general}.  For either cutoff loss $L=L_i$, put
$H=\log(en/k)$.  Expanding
$e^{uHL}=(1+(e^{uH}-1))^L$ shows that its expectation is at most one
plus a geometric series with ratio
$C\sqrt{x}((e/x)^u-1)$, where $x=k/n$.  This is $O(u)$ uniformly on
$0<x\le1$ for $u\le1/4$, proving~\eqref{rec:loss-mgf}.

For~\eqref{rec:loss-tail}, first restrict to $a\le.8m$ and $r<m/16$.
The lower removal factor is at least $.1m$, so on $a-r\ge x$,
\[
 \Pr(a\le.8m,r<m/16,a-r\ge x)
 \le\frac{(C\sqrt{k/m})^j}{\binom{x}{j}}
 \quad(1\le j\le x).
\]
For small $k/n$ take $j=x$.  On the remaining density range $H$ is
bounded; take $j=\lfloor x/M\rfloor$ with a sufficiently large absolute
$M$, using ${x\choose j}\ge(x/j)^j$.  The bounded $x<M$ cases are
absorbed in the constant.  If $a>.8m$ and $r<m/16$, then $k>.6n$ is
impossible by $r\ge k-m$.  For $k\le.6n$ the event requires a deviation
of at least $.2m$, while also $k>.8m$.  Its probability is $Ce^{-ck}$
by~\eqref{rec:split-tail}; here $H$ is bounded and $x\le k$.
\end{proof}

\subsection{Three ways to close the level induction}
We have controlled both the imbalance of the two tuple occupancies and
the loss from an occupancy to a child representation level.  The next step
is to turn those two controls into the first bound in
Theorem~\ref{rec:main-harmonic-cycle}.  Three potentials do this with
different sources of strictness: a fixed multiplicative margin, a large
square-root correction, or a unit square-root correction coupled to a
fixed-level estimate.  We first complete the constant-margin proof.
The two square-root alternatives then show how to replace that margin;
the last also supplies the fixed-level estimate
\eqref{rec:fixed-level-bootstrap}.

\medskip\noindent\textbf{Capped potential comparisons.}\enspace
The three inductions use the same elementary estimates.  For $c>0$ set
\[
 E_c(x)=\begin{cases}x\log(ec/x),&0<x\le c,\\c,&x>c,\\0,&x=0,
 \end{cases}
 \qquad R_c(x)=\sqrt{\min(x,c)}.
\]
Take $c=\alpha m$, where $0<\alpha\le1/8$ is fixed, and write
$H=\log(en/k)$, $D=(a-k/2)^2/k$ and
$U_1=(a-r)\mathbf1_{\{r<c\}}$, with $U_2$ defined in the other half.
The identities $E_{2c}(k)=2E_c(k/2)$ and $a+b=k$ give
\[
 0\le E_{2c}(k)-E_c(a)-E_c(b)
 \le a\log(2a/k)+b\log(2b/k)\le4D.
\]
For the first inequality on the right, the magnitude of the negative
second derivative of $E_c$ is at most $1/x$; integrating from the
balanced split compares its gap with the displayed entropy gap.
The last inequality is the relative-entropy bound by the corresponding
chi-square divergence on two points.  Concavity and $E_c(0)=0$ also give
\[
 E_c(a)-E_c(r)
 \le U_1\{1+\log_+(c/a)\}.
\]
When $a\ge k/4$, the braces are at most $H$.  When $a<k/4$, their
excess over $H$ costs at most $a\log(k/a)\le k/e$, while
$D\ge k/16$.  Thus
\begin{equation}\label{rec:capped-entropy-comparison}
 E_{2c}(k)-E_c(r)-E_c(s)\le C D+H(U_1+U_2).
\end{equation}
All these inequalities include zero occupancies by continuity.

The map $x\mapsto\sqrt{\min(x,c)}$ is a contraction as a function
of $\sqrt x$.  Hence
$|R_c(a)-R_c(k/2)|\le\sqrt{2D}$, and
$R_c(a)-R_c(r)\le\sqrt{U_1}$.  With $s_0=\min(k,2c)$ this yields
\begin{equation}\label{rec:capped-root-comparison}
 R_{2c}(k)-R_c(r)-R_c(s)
 \le-(\sqrt2-1)\sqrt{s_0}+2\sqrt{2D}
                              +\sqrt{U_1}+\sqrt{U_2}.
\end{equation}
These comparisons separate occupancy imbalance from the loss of
representation level, so the estimates in Lemma~\ref{rec:removal-splitting}
can be applied without conditioning on either event.

\medskip\noindent\textbf{A constant multiplicative margin.}\enspace

First take $x_*=1/100$ and define
\[
 H_N(j)=N h(j/N),\qquad
 h(y)=\begin{cases}y(1+\log(x_*/y)),&0<y\le x_*,\\
 x_*,&y>x_*,\\0,&y=0.\end{cases}
\]
The function is nondecreasing and concave, and
$H_n(k)\ge c\Lambda_n(k)$.  Put
$\Delta=H_n(k)-H_m(r)-H_m(s)\ge0$.
There are absolute $\beta,C>0$ such that
\begin{equation}\label{rec:constant-margin-mgf}
 \E e^{\beta\Delta}\le C.
\end{equation}
Here is the loss calculation.  For $k/n\ge3/5$, the diagram containment
bound gives $r,s\ge m/5$ and the gap is zero.  Otherwise introduce
$H_m(a)+H_m(b)$ and use the preceding entropy comparison.
If $a>4m/5$ but $r<x_*m$, containment implies
$k<(1+x_*)m$.  Therefore $D\ge c m$, which bounds the entire capped
loss, at most $x_*m$.  For all other occupancies the same calculation,
with $1+\log_+(4x_*m/k)$ in place of $H$, gives
\[
 \Delta\le CD+J(V_1+V_2),\quad D=(a-k/2)^2/k,
 \quad J=1+\log_+(4x_*m/k),
\]
where $V_1=(a-r)\mathbf1_{\{r<x_*m,a\le4m/5\}}$, and similarly for
$V_2$.  On this indicator $m-a-r\ge m/10$.
The binomial moments in~\eqref{rec:loss-binomial-general} imply
\[
 \E e^{tJV_i}\le1+\sum_{l\ge1}
       [(e^{tJ}-1)C_{1/10}\sqrt{k/n}]^l.
\]
The ratio is uniformly small for a sufficiently small fixed $t$.
H\"older and~\eqref{rec:split-tail} prove
\eqref{rec:constant-margin-mgf}.

For all sufficiently large $n$, uniformly over $k\ge2$, the same
component law satisfies
\begin{equation}\label{rec:constant-margin-events}
 \Pr(r,s\ge1)\ge1/3,\qquad \Pr(r=s=0)\le1/12.
\end{equation}
For density at least $3/5$ this is deterministic.  Otherwise, when $k$
is large the two occupancies lie between $k/4$ and $4m/5$ except with
vanishing probability, and a zero child level would make a $V_i$ at least
$k/4$.  The preceding exponential estimate excludes this.  For bounded
$k\ge2$, the probability of any removal loss tends uniformly to zero by
the $l=1$ bound.  The binomial mixture has two positive occupancies with
probability at least $1-2^{1-k}\ge1/2$, and cannot have both zero.

Choose a finite base threshold for~\eqref{rec:constant-margin-events}.
Every nontrivial sweep norm at those finitely many sizes is strictly below
one: equality in a product of orthogonal projections would give a vector
fixed by every cube-edge transposition, hence by $S_n$.  Choose
$0<A\le1/2$ so the finite base norms are at most $e^{-2A}$.
We induct with
\begin{equation}\label{rec:constant-margin-invariant}
 W_n(\lambda)\le e^{-A-\varepsilon H_n(k)}\quad(k\ge2).
\end{equation}
Level one is annihilated by exact one-card uniformity.  Let
$G=\{r\ne1,s\ne1\}$ and $b=\#\{r,s\text{ that are at least }2\}$.
After dividing the recursive bound by the proposed right side, the
remaining factor is
\[
 \E[\mathbf1_G e^{-A(b-1)}e^{\varepsilon\Delta}].
\]
Without the last factor this is bounded uniformly below one: the event
of two positive levels loses at least $1-e^{-A}$, whereas only two zero
levels can gain $e^A-1\le2(1-e^{-A})$.
Equations~\eqref{rec:constant-margin-events} bound the factor by
$1-(1-e^{-A})/6$.  Moreover
\eqref{rec:constant-margin-mgf} and H\"older give
$\E e^{\varepsilon\Delta}\le C^{\varepsilon/\beta}=1+O(\varepsilon)$.
Since $\mathbf1_Ge^{-A(b-1)}\le e^A$, reinstating
$e^{\varepsilon\Delta}$ costs only $O(\varepsilon)$, uniformly.
Also choose $\varepsilon H_n(k)\le A$
at the base sizes.  This proves~\eqref{rec:level-contraction}.

\medskip\noindent\textbf{A large square-root correction.}\enspace
A second invariant dispenses with the constant margin by inserting a
square-root term.  With $\alpha=1/16$, set
\[
 G_N(j)=N\phi(j/N),\qquad
 \phi(y)=\begin{cases}y\log(e\alpha/y),&0<y\le\alpha,\\
 \alpha,&y\ge\alpha,\\0,&y=0,\end{cases}
\]
\[
 F_N(j)=G_N(j)+100\sqrt{\min(j,\alpha N)}.
\]
For $Y=F_n(k)-F_m(r)-F_m(s)$ and $s_0=\min(k,\alpha n)$,
\begin{equation}\label{rec:sqrt-large-drift}
 Y\le-100(\sqrt2-1)\sqrt{s_0}
       +C_{100}\{1+D+H(U_1+U_2)\},
\end{equation}
where $H=\log(en/k)$, $D=(a-k/2)^2/k$, and the $U_i$ are the
uncut losses in~\eqref{rec:loss-tail}.
Indeed~\eqref{rec:capped-entropy-comparison} and
\eqref{rec:capped-root-comparison}, together with
$\sqrt D\le1+D$ and $\sqrt{U_i}\le U_i$ for integer losses, give
\eqref{rec:sqrt-large-drift}.  The split and loss tails imply, for small
$\tau\ge0$, a bound
\[
 \E e^{\tau C_{100}\{1+D+H(U_1+U_2)\}}\le e^{C'\tau}.
\]
For example the tail in~\eqref{rec:loss-tail} bounds
$\E e^{cHU_i}$ uniformly for a small fixed $c>0$; H\"older and
convexity then give the displayed bound at smaller exponents.
Since $s_0\ge\alpha k$, one fixed threshold $K$ makes
$\E e^{\tau Y}\le1$ for every $k>K$ and all sufficiently small $\tau$.

For $2\le k\le K$, take $n$ large enough that the potentials are
uncapped at these levels.  Put $L=(a-r)+(b-s)$, and decompose the
gap without conditioning on the loss event:
\[
 Y=Q-100S+R,\qquad
 Q=a\log(2a/k)+b\log(2b/k),\quad
 S=\sqrt a+\sqrt b-\sqrt k,
\]
where $0\le R\le(H+C_K)L$.  The binomial mixture gives
$\E Q\le4\E D\le1$.  Both occupancies are positive with probability
at least $1/2$, and then $S\ge1/2$, so $\E S\ge1/4$.
The first removal moment gives $\E L\le C_K e^{-H/2}$.
Consequently $\E Y\le-24+o(1)\le-12$ for all sufficiently large $n$,
uniformly over these finitely many levels.

The same bound also justifies a uniform Taylor remainder.  We have
$|Y|\le C_K+KH\mathbf1_{\{L>0\}}$ and
$\Pr(L>0)\le C_K e^{-H/2}$, and therefore
\[
 M_K:=\sup\E\big[Y^2e^{|Y|/(4K)}\big]<\infty,
\]
where the supremum is over these levels, sizes, and component laws.
For $0<\tau\le\min\{1/(4K),12/M_K\}$, Taylor's inequality gives
$\E e^{\tau Y}\le1+\tau\E Y+\tau^2 M_K/2\le1-6\tau$.
The remaining finitely many sizes are covered by shrinking $\tau$ using
the strict norm gap.  Induction now proves the invariant
$W_n(\lambda)\le e^{-\tau F_n(k)}$.

\medskip\noindent\textbf{A unit correction and fixed-level bootstrap.}\enspace
A third version uses coefficient one instead of $100$.  Use the
preceding definition of $G_N$ with $\alpha=1/8$, and put
$\Phi_N(j)=G_N(j)+\sqrt{\min(j,\alpha N)}$.
The cut losses $L_i$ in Lemma~\ref{rec:removal-splitting} give
\begin{equation}\label{rec:sqrt-unit-drift}
 \Phi_n(k)-\Phi_m(r)-\Phi_m(s)
 \le-c_0\sqrt k+C_0(a-b)^2/k+C_0H(L_1+L_2).
\end{equation}
Here are the capped cases in this comparison.  Put $c=\alpha m$.
If $a,b\in[.4k,.6k]$ and $c/k\le.4$, both child square roots are
capped, so their saving over the root is
$(2-\sqrt2)\sqrt c\ge(2-\sqrt2)\sqrt k/4>.1\sqrt k$;
we used $c/k\ge1/16$.  If $c/k\ge.4$, each child square root
is at least $\sqrt{.4k}$ and the root is at most $\sqrt k$, giving
the same saving.  Outside this occupancy interval, $D\ge k/100$,
and the square-root gap, at most $\sqrt k$, is bounded by
$-.1\sqrt k+110D$.  This proves the required no-loss bound.

The entropy comparison~\eqref{rec:capped-entropy-comparison} and the
root loss $\sqrt{L_i}\le L_i\le HL_i$ cover the cut losses.
The only omitted case has $a>3m/4$ and $r<m/8$, or its counterpart.
Containment gives $k<9m/8$, hence $D>m/32$; the entire capped
entropy and square-root loss is at most $m$, and so is bounded by
$32D$.  These bounds prove~\eqref{rec:sqrt-unit-drift} with
$c_0=1/10$ and an absolute $C_0$.
Equations~\eqref{rec:split-tail} and~\eqref{rec:loss-mgf} imply
\[
 \E e^{\tau(\Phi_n(k)-\Phi_m(r)-\Phi_m(s))}
       \le e^{\tau(B-c_0\sqrt k)}
\]
for fixed $B$ and all sufficiently small $\tau$.
One may start this induction through the additional quantitative fact
\begin{equation}\label{rec:fixed-level-bootstrap}
 \max_{\lambda:n-\lambda_1=k}W_n(\lambda)=O_k(n^{-1/4})
 \qquad(k\ge1\text{ fixed}).
\end{equation}
Induct on $k$.  A positive removal loss has probability $O_k(n^{-1/2})$.
Without losses, two positive child levels are smaller than $k$ and obey
the induction hypothesis.  All labels in one half has probability at most
$2^{1-k}$, giving for $k\ge2$
$M_k(2m)\le2^{1-k}M_k(m)+O_k(m^{-1/4})$.
Since $2^{1-k}2^{1/4}<1$, this proves
\eqref{rec:fixed-level-bootstrap}; level one vanishes.
For any fixed upper level $K$, the estimate is eventually stronger
than $e^{-\tau\Phi_n(k)}$ simultaneously for $1\le k\le K$ if
$\tau\le1/(8K)$, since $\Phi_n(k)=k\log n+O_K(1)$ there.
Shrinking $\tau$ handles the remaining finite sizes.  Thus it supplies
one common small $\tau$ for all bounded levels, and then
\eqref{rec:sqrt-unit-drift} closes induction for the larger levels.
Thus all three invariants prove the first assertion of
Theorem~\ref{rec:main-harmonic-cycle}, while
\eqref{rec:fixed-level-bootstrap} records a separate fixed-level estimate.

\subsection{Cycle moments from positive tensor traces}
We next prove the density estimate.  In a node of size $2m$ in the
recursive network, form the bipartite graph of input and output switch
cells, with one edge for each tracked card.  It is a union of paths and
even cycles, including double edges.  Let $c_v^*$ be its cycle count;
let $c_v$ be the count when all cards are tracked.  A node of size $2^j$ has height $j$.  Write $C_j^*,C_j$
for their sums over all nodes at height $j$.  Then
$C_j^*\le k/2$ and $C_j^*\le C_j$.
The full count at a node is the number of cycles of the relative
permutation of its two children.  In particular it depends only on
the randomness below that node, irrespective of the tracks arriving
from above.

For $q\ge1$ put
$Z_l(q)=\E_{g\sim K_{2^l}}q^{\#\mathrm{cycles}(g)}$.
This scalar generating function is defined for every real $q\ge1$.
For an integer $q$ it is also the unnormalized trace of $K_{2^l}$ on
$(\mathbb C^q)^{\otimes2^l}$.  The tensor identity in
Lemma~\ref{cas:swap-recursion}, specialized to this integer alphabet, gives
\begin{equation}\label{rec:cycle-trace-recursion}
 Z_l(q)\le Z_{l-1}(q)Z_{l-1}(1+(q-1)/2)
 \qquad(q\in\mathbb Z_{\ge1}).
\end{equation}
For completeness, divide the positive child operator by its trace and
call the resulting density matrix $D$.  Expanding the crossing average
over a uniformly chosen subset $S$ of tensor positions gives
$Z_l(q)=Z_{l-1}(q)^2\E_S\operatorname{tr}(D_S^2)$.
The partial trace of a permutation is its induced permutation on $S$,
multiplied by $q$ for each cycle avoiding $S$.  Therefore its norm is
at most that scalar.  Bounding $\operatorname{tr}(D_S^2)$ by
$\|D_S\|_{\op}$ and averaging $S$ contributes at most
$1+(q-1)2^{-|C|}\le1+(q-1)/2$ per cycle, proving the recursion.

Two useful explicit versions are
\begin{align}
 2^{-l}\log Z_l(1+2^h)&\le2^h(.81)^l
             &&(h=0,1,\ldots),\label{rec:cycle-081}\\
 \log Z_l(1+w)&\le w(2^l)^\gamma,\qquad
 \gamma=\log_2(1+\log_2(3/2))<1
             &&(w=2^h,\ h\in\mathbb Z_{\ge0}).\label{rec:cycle-gamma}
\end{align}
For $q>2$ the normalized induction factor is $3/4$; for $q=2$ use
$Z_l(3/2)\le Z_l(2)^{\log_2(3/2)}$, giving factor
$(1+\log_2(3/2))/2<.81$.  The same two cases yield the sharper
exponent in~\eqref{rec:cycle-gamma}, since $3/2\le2^\gamma$.
The base $Z_0(q)=q$ satisfies both bounds.

Here is the conditional form needed to sum heights.  Let $\mathcal F_h$ be the sigma-field generated by all switches
at heights at most $h$, with the trivial sigma-field when $h<1$.
Reveal $\mathcal F_{j-l-1}$.  Each child of a height-$j$ node has conditional
average $A^*UA$, with $U$ a fixed unitary and $A$ a tensor product of
$l$-coordinate sweeps on $2^l$-slot cells.  Independence of the two
children and Hilbert--Schmidt Cauchy--Schwarz show
\begin{equation}\label{rec:conditional-cycle-trace}
 \E[q^{C_j}\mid\mathcal F_{j-l-1}]
      \le Z_l(q)^{n/(2\cdot2^l)}.
\end{equation}
Indeed, writing $F_0,F_1$ for the two conditional tensor mean
operators, $\operatorname{tr}(F_0^*F_1)\le\operatorname{tr}((AA^*)^2)
\le\operatorname{tr}(AA^*)$, and the trace factors over the cells.
Fresh switches in different height-$j$ nodes are conditionally independent.

Choose $\delta>0$ with $4\delta<1-\gamma$.  Partition heights into
bands $[b,2b-1]$ for powers of two $b$, and separate alternating bands.
For $b\ge2$, take $l=b/2$ and round $2^{2\delta b}$ up to an integer
alphabet, then its excess above one up to a power of two.  Equations
\eqref{rec:cycle-gamma}--\eqref{rec:conditional-cycle-trace} give
\[
 \E[2^{2\delta bC_j}\mid\mathcal F_{b/2-1}]
 \le\exp(Cn2^{-\eta b}),\qquad
 \eta=(1-\gamma)/2-2\delta>0.
\]
Conditional H\"older over the at most $b$ heights gives the same bound
for $2^{2\delta\sum C_j}$ in the band.  A smaller band of the same
parity ends below $b/2$, so its factor is measurable for this conditioning.
Integrate the highest band first and iterate; Cauchy--Schwarz combines
the two parities.  For $J\ge2$ the result is
\begin{equation}\label{rec:high-cycle-moment}
 \E2^{\delta\sum_{j\ge J}C_j}
       \le\exp(Cn2^{-\eta'J}),\qquad \eta'=\eta/2.
\end{equation}
The use of full counts is essential: replace tracked counts by full
counts in every band where conditional integration will be used before
performing that integration.  Tracked lower counts can otherwise depend
on the routes chosen above them.

Taking $J$ a sufficiently large constant multiple of $\log_2(en/k)$,
the deterministic bound on lower heights and
\eqref{rec:high-cycle-moment} give, uniformly over initial tuples,
\begin{equation}\label{rec:tracked-cycle-moment}
 \E2^{\delta\sum_jC_j^*}\le e^{C\Lambda_n(k)}.
\end{equation}
Equivalently one may use~\eqref{rec:cycle-081}: in a band starting at
$b$, the conditional logarithmic bound is
$(n/2)2^{\lceil2\delta b\rceil}(.81)^{\lfloor b/2\rfloor}
\le Cn e^{-\gamma_0b}$ for a sufficiently small fixed $\delta$.
Using tracked deterministic bounds below
$\gamma_0^{-1}\log(n/k)$ and full counts above it proves the same
estimate with explicit geometric decay in band height.

\subsection{An entropy coding proof of the cycle moment}
The tensor argument proves the required moment.  The following coding
argument also controls the cycle count under an arbitrary change of
the switch law, charging its increase to relative entropy.  This is
the additional conclusion we retain from this alternative proof.
In this subsection entropy is measured in bits.

At height $t$ the switch array consists of both outer layers of every
node of size $2^t$; all these raw switch bits are independent under the
reference fair law.  Let $Q$ be any probability law on the full array
of raw switches at all heights of the $n$-position reflected network.
Write $\mathcal B_t$ for the vector of bits at height $t$ and define the
nonnegative number
\[
 I_t=|\mathcal B_t|-H_Q(\mathcal B_t\mid\mathcal B_1,\ldots,
 \mathcal B_{t-1}),
\]
where conditional Shannon entropy is already averaged under $Q$.  Thus $I_t\ge0$ and $\sum_tI_t=D(Q\Vert\mathrm{fair})$.
\begin{proposition}[Cycle costs under a change of law]
\label{rec:entropy-deficit}
For any fixed input tuple of $k$ distinct positions and any law $Q$
on the raw switches, the actual tracked cycle counts satisfy
\begin{equation}\label{rec:cycle-tilted-entropy}
 \E_Q\sum_j C_j^*\le C\Lambda_n(k)+C D(Q\Vert\mathrm{fair}),
\end{equation}
with an absolute constant $C$.
\end{proposition}
\begin{proof}
With $h=\lfloor j/2\rfloor$, the full cycle counts obey
\begin{equation}\label{rec:cycle-window-entropy}
 \E_Q C_j\le Cne^{-cj}+\frac{C}{j}
                         \sum_{t=j-h}^{j-1}I_t\qquad(j\ge2).
\end{equation}
To prove this, code the window bits given smaller heights.  At one
height-$j$ node, with child size $m=2^{j-1}$, send literally all window
bits except the first $h$ input layers of one child.  The relative child
permutation is now $D_0A$, with $D_0$ known and $A$ a butterfly with
$hm/2$ unknown bits.  If $D_0A$ has $r$ cycles, send $r$, one
representative per cycle, and the ordered list of cycle lengths.  These
data cost at most
$C[\log_2(2m)+r\log_2(em/r)]$ bits, using
${m\choose r}$ and ${m-1\choose r-1}$.

Traverse each cycle from its representative.  Send each as-yet-unknown
switch bit used on its $A$ paths, except on the last path of that cycle:
its required output is already known, so its unique butterfly route
recovers all those bits without sending them.  The specified lengths tell the decoder exactly when this omission occurs.
For each fixed choice of the public orders described next, use a fixed-width
field for $r$, an enumerative code for its representative set, and an
enumerative code for the composition of $m$ into the transmitted lengths.
The decoder then reads a bit precisely when the traversal first queries
an unknown switch outside a closing path.  It consequently knows when the
whole message has ended.  The resulting code is prefix-decodable.  Every
switch is recovered after all positions have been traversed.

Use independent public random priority orders on vertices to choose a
uniform representative within each cycle and then to order the transmitted
representative set.  The decoder uses that same second priority order;
no extra factorial cost is needed to specify the cycle order.  For any set of $s$ butterfly
paths, the number of internal contacts, counted at both ends, is at most
$s\log_2s$.  Indeed, splitting a block into two predecessor blocks of
sizes $s_1,s_2$ adds at most $2\min(s_1,s_2)$ contacts, and
\[
 s_1\log_2s_1+s_2\log_2s_2+2\min(s_1,s_2)
           \le(s_1+s_2)\log_2(s_1+s_2).
\]
The closing path of a length-$s$ cycle therefore has, on average, at
least $h-\log_2s$ contacts outside its cycle.  Each outside cycle is
later with probability $1/2$, so those switches are then unknown.
The expected saving is at least
\[
 \frac{r}{2}\{h-\log_2(m/r)\}
   -C[\log_2(2m)+r\log_2(em/r)].
\]
This remains a valid lower bound if some terms are negative, and covers
cycles of length one.  The public orders are independent of the switch array.  Conditional on
the smaller heights and on those orders, the code remains a prefix code
for the entire window: different height-$j$ nodes use disjoint sets of its
bits.  Its entropy is at most its expected length.  Average over the public
orders and the smaller heights.  The difference between the literal
window length and its conditional entropy is exactly
$\sum_{t=j-h}^{j-1}I_t$, by the entropy chain rule.  Thus the saving just
computed is charged to that window deficit, without assuming independence
of the bits under $Q$.
The inequality
$Cr\log_2(em/r)\le(h/4)r+C'me^{-c'h}$
and the $n/2^j$ nodes give~\eqref{rec:cycle-window-entropy}.

At each height the tracked count is at most $k$ and at most the full
count.  Moreover each $I_t$ occurs with bounded total coefficient in
the windows, since $\sum_{j>t,\ j-\lfloor j/2\rfloor\le t}1/j=O(1)$.
Sum the deterministic tracked bound below a sufficiently large multiple
of $\log(en/k)$ and~\eqref{rec:cycle-window-entropy} above it.
The first sum and the exponentially decaying errors cost
$O(\Lambda_n(k))$; the bounded window coefficients cost
$O(\sum_t I_t)$.  This proves~\eqref{rec:cycle-tilted-entropy}.
\end{proof}
Apply this to the law tilted by $2^{\varepsilon\sum C_j^*}$.
The variational identity
\[
 \log_2\E2^{\varepsilon\sum C_j^*}
 =\varepsilon\E_Q\sum C_j^*-D(Q\Vert\mathrm{fair})
\]
proves~\eqref{rec:tracked-cycle-moment} when $\varepsilon C<1$.
Thus the tensor-trace argument and the coding argument each supply a
complete, separate proof of the required cycle moment.

\subsection{From routed cycles to densities and multiplicities}
Consider specified injective endpoints for $h$ links in a size-$2m$
node.  Let $e$ be the total number of shared input and output pairs,
and let $c$ be the cycle count in their union.  The links admit exactly
$2^{h-e+c}$ proper child colorings.  Each has switch probability
$2^{-2h+e}$.  If $R=(2m)^h p$ is the provisional transition density,
the child recursion is
\begin{equation}\label{rec:routing-density-recursion}
 R=\sum_{\text{proper colors}}2^{-h+e}R_0R_1
     =2^c\operatorname{Avg}_{\text{proper colors}}R_0R_1.
\end{equation}
Jensen with exponent $1+\delta$ charges a factor $2^{\delta c}$.
After unrolling to one-slot leaves,
\[
 (n^kp(x,y))^{1+\delta}
 \le n^k\E_{K_n}[\mathbf1_{\{\pi x=y\}}
                         2^{\delta\sum C_j^*}].
\]
The remaining coefficients are the actual switch probabilities: each
link is counted once per height, and specified recursive routings are
disjoint events.  Summing in $y$ and using
$(n)_k/n^k\le1$ proves~\eqref{rec:tuple-moment}.
Symmetry supplies the corresponding column estimate.  We will also use
the provisional normalization directly:
\begin{equation}\label{rec:provisional-moment}
 \sup_x\sum_y n^{-k}(n^kp(x,y))^{1+\delta}
 \le e^{C\Lambda_n(k)}.
\end{equation}
Indeed summing the preceding unrolled inequality proves this before the
factor $(n)_k/n^k$ is inserted.  Equivalently it follows from the density
bound after enlarging $C$, since
\[
 \log\frac{n^k}{(n)_k}
 =\sum_{i=0}^{k-1}-\log(1-i/n)
 \le\int_0^k-\log(1-x/n)\,dx\le k.
\]
The integral is finite also at $k=n$.

The row and column bounds imply bounded operators
$K_n:L^1\to L^{1+\delta}$ and
$K_n:L^{(1+\delta)/\delta}\to L^\infty$, each with norm at most
$M=e^{C\Lambda_n(k)}$ after enlarging $C$.
Riesz--Thorin interpolation decreases the reciprocal exponent by
$\delta/(1+\delta)$ at each application.  After
$p_0=\lceil(1+\delta)/\delta\rceil$ applications, with probability-space
inclusion for the last step, $K_n^{p_0}:L^1\to L^\infty$ has norm
at most $M^{p_0}$.  Equivalently,
\begin{equation}\label{rec:tuple-pointwise-smoothing}
 \sup_{x,y}|X_{n,k}|K_n^{p_0}(x,y)\le e^{C\Lambda_n(k)}.
\end{equation}
The diagonal kernel bound gives
\eqref{rec:tuple-flat-trace}.

There are two ways to use the label multiplicity $f_\beta$.
For the first transfer, the goal is to use the $f_\beta$ equivalent
slot copies of one position vector.  Clip the entries $p(x,y)$ of $K_n$ above
$n^{-k}\sqrt{f_\beta}$, setting them to zero, and call the retained
matrix $M_0$.  The provisional moment bound~\eqref{rec:provisional-moment} and
symmetry give
\begin{equation}\label{rec:tail-clipping-error}
 \|K_n-M_0\|_{\op}\le e^{C\Lambda_n(k)}f_\beta^{-\delta/2}.
\end{equation}
Since every row sum is at most one and $|X_{n,k}|\le n^k$,
$\|M_0\|_{\HS}^2\le\sqrt{f_\beta}$.
Clipping commutes with relabeling the tuple slots.  Compressing to
position type $\lambda$, whose label type is $\beta$, therefore repeats
each singular value at least $f_\beta$ times.  Its compressed norm is
at most $f_\beta^{-1/4}$.  If $\log f_\beta$ exceeds a sufficiently
large constant times $\Lambda_n(k)$, these two estimates give
$W_n(\lambda)\le e^{-c\log f_\beta}$; otherwise
\eqref{rec:level-contraction} applies.  This proves
\eqref{rec:combined-tail-norm}.

The clipping proof now gives the combined level-and-tail estimate.  To
obtain the explicit tail bound~\eqref{rec:tail-contraction} by a different
mechanism, average a matrix coefficient over all slot copies.  Tensor a maximizing unit vector in $V_\lambda$ with an
orthonormal basis of $V_\beta$, obtaining orthonormal functions $u_j$,
$1\le j\le f_\beta$, in probability norm.  Define
\[
 F(x,y)=\frac{1}{f_\beta}\sum_j\overline{u_j(x)}u_j(y).
\]
Then $\E(F\mathcal K)=W_n(\lambda)$ and
$\E|F|^2=f_\beta^{-1}$.  The diagonal
$G(x)=f_\beta^{-1}\sum_j|u_j(x)|^2$ is slot-invariant and has mean
one.  Each orbit of the slot-permutation group is the set of $k!$
orderings of one $k$-element set, so has probability
$\binom nk^{-1}$.  Nonnegativity therefore gives
$G(x)\le\binom nk$.  Cauchy--Schwarz gives the same bound
for $|F(x,y)|$.  With $p=1+\delta$ and $p'=p/(p-1)>2$, H\"older gives
\begin{equation}\label{rec:tail-averaged-coefficient}
 W_n(\lambda)\le\|\mathcal K\|_p\|F\|_{p'}
 \le e^{C\Lambda_n(k)}{n\choose k}^{1-2/p'}f_\beta^{-1/p'}.
\end{equation}
This proves~\eqref{rec:tail-contraction} with $\zeta=1/p'$.
Since $D_\lambda\le{n\choose k}f_\beta$, geometric interpolation with
\eqref{rec:level-contraction} proves
\eqref{rec:combined-dimension-norm}.

Finally, the tuple trace gives a completion that needs no tail-dimension
interpolation.  At exact level $k$ the multiplicity in $X_{n,k}$ is
$f_\beta$, while the regular multiplicity is
$D_\lambda\le{n\choose k}f_\beta$.  For an integer $p>p_0$, positivity and the
estimates \eqref{rec:level-contraction} and \eqref{rec:tuple-flat-trace} yield
\begin{equation}\label{rec:tuple-regular-completion}
 \operatorname{tr}_{\rm reg}K_n^p-1
 \le\sum_{k=1}^{n-1}
 e^{-[(p-p_0)c-C-1]\Lambda_n(k)}.
\end{equation}
Alternatively~\eqref{rec:combined-dimension-norm} bounds the regular
Fourier sum directly by $\sum_{k\ge1}2^k e^{-c_1M\Lambda_n(k)}$ once
$c_2M\ge2$.  Both sums tend to zero for a sufficiently large fixed
exponent.  Schatten H\"older gives
$\|T_n^p|_{\mathbf1^\perp}\|_{\HS}^2\le\operatorname{tr}_{\rm reg}K_n^p-1$ without
assuming normality of $T_n$.  Thus these arguments mix the entire
permutation law after an absolute number of sweeps, each consisting
of exactly $d$ physical shuffles.
